% Capitolo 1: Introduzione

La polarizzazione della luce descrive l'orientazione e il comportamento temporale del campo elettrico nel piano trasversale alla direzione di propagazione.A differenza dell’intensità e della frequenza, la polarizzazione non è direttamente percepibile, ma codifica informazioni essenziali sulle proprietà del mezzo attraversato: struttura cristallina, stato di sollecitazione meccanica, chiralità molecolare. La polarimetria, ossia la misura quantitativa dello stato di polarizzazione, è per questa ragione uno strumento trasversale a ottica, scienza dei materiali, chimica analitica e ingegneria strutturale, con applicazioni in fotoelasticità, diagnostica biomedica, controllo qualità industriale e osservazione astronomica.


\section{Dalla misura puntuale all'immagine}
\label{sec:polarimetria_classica}
L'approccio tradizionale alla polarimetria impiega un rivelatore puntuale (tipicamente un fotodiodo) posto dietro un sistema di analizzatore e lamine ritardatrici. Ruotando l'analizzatore si ricostruiscono i quattro parametri di Stokes del fascio incidente: configurazione descritta, ad esempio, nel manuale di laboratorio del corso di Fisica Sperimentale del Politecnico di Milano~\cite{manuale_polarizzazione}, preso come punto di partenza per questa tesi. Il limite principale di questo schema è la sua natura \emph{zero-dimensionale}: il rivelatore fornisce una singola misura scalare, integrata sull'intera sezione del fascio. Se il campione presenta variazioni spaziali delle proprietà ottiche, come accade nella maggior parte delle situazioni di interesse pratico, l'informazione spaziale viene irrimediabilmente perduta. Un secondo limite, per quanto meno stringente, è economico: un polarimetro di Stokes puntuale di grado commerciale, in cui un analizzatore o una lamina ritardatrice motorizzata ruota davanti a un fotodiodo, ha un costo dell'ordine di alcune migliaia di euro.

L'idea di sostituire il rivelatore puntuale con un sensore a matrice trasforma il polarimetro da strumento scalare a strumento di imaging: ogni pixel diventa un polarimetro indipendente, e la misura restituisce una mappa bidimensionale dello stato di polarizzazione. Le soluzioni commerciali per questo scopo si distribuiscono in due famiglie con prestazioni e costi assai diversi. Le telecamere a divisione del piano focale, costruite attorno a un array di micro-polarizzatori integrato direttamente sul sensore CMOS, forniscono mappe di polarizzazione in un singolo scatto a un costo di alcune migliaia di euro, ma misurano la sola polarizzazione lineare ($S_0$, $S_1$, $S_2$): l'assenza di elementi ritardatori sul piano focale impedisce loro l'accesso alla componente circolare $S_3$. La misura del vettore di Stokes completo su un campo bidimensionale richiede invece sistemi attivi, con modulatori a cristalli liquidi o compensatori rotanti motorizzati. % TODO: vero? questi strumenti esistono davvero? poi frase sembra implicare due famiglie ma la seconda famiglia sembra introdotta con una frase a parte

% TODO: rifare paragrafo sotto, prima dicono che sono a basso costo poi che costano migliaia di euro, inoltre andrebbe verificato fulltext delle fonti, le allego in biblio
Strumenti polarimetrici a basso costo basati su smartphone sono già documentati in letteratura per la spettropolarimetria atmosferica~\cite{Burggraaff2020ispex}, per la diagnostica biomedica~\cite{Gonzalez2020colposcope,Louie2021skin} e nella didattica della fisica~\cite{Gallitto2024malus,Bernard2020polarimeter}, mentre la ricostruzione dell'intero vettore di Stokes a immagine in un singolo scatto è oggetto di ricerca recente~\cite{Baek2022lensless,Gu2022fullstokes}. Un prodotto commerciale capace di quest'ultima misura ha un costo dell'ordine delle decine di migliaia di euro e resta accessibile soprattutto ai laboratori ben attrezzati.

% TODO: trovo riduttivo parlare di studenti, questo è un lavoro utle in generale
La diffusione capillare degli smartphone, con sensori CMOS ad alta risoluzione e acquisizione RAW, suggerisce una possibilità concreta: impiegare come rivelatore un dispositivo già posseduto dalla quasi totalità degli studenti per ottenere mappe del vettore di Stokes completo a un costo dell'ordine delle centinaia di euro, due ordini di grandezza inferiore a quello dei sistemi commerciali capaci della stessa misura.


\section{Obiettivo del lavoro}
\label{sec:obiettivo}

Questa tesi si propone di progettare, realizzare e validare un polarimetro a immagine bidimensionale a basso costo, utilizzando come rivelatore il sensore CMOS di uno smartphone commerciale. Il sistema deve ricostruire il vettore di Stokes completo pixel per pixel e generare mappe spazialmente risolte del grado di polarizzazione, dell'angolo di polarizzazione lineare, della ritardanza e dell'orientazione dell'asse veloce. Per validarne le capacità, vengono analizzati campioni rappresentativi di birifrangenza, attività ottica e fotoelasticità.

Il \cref{ch:teoria} richiama i fondamenti teorici necessari alla comprensione del metodo di misura. Il \cref{ch:apparato} descrive l'apparato sperimentale e il protocollo di acquisizione. Il \cref{ch:campioni} presenta i campioni selezionati. Il \cref{ch:analisi} documenta l'analisi dei dati. Il \cref{ch:risultati} discute i risultati e i limiti del metodo. Il \cref{ch:conclusioni} riassume le conclusioni e indica possibili sviluppi futuri.
